\documentclass[12pt,a4paper]{article}
\usepackage[utf8]{inputenc}
\usepackage[T1]{fontenc}
\usepackage[polish]{babel}
\usepackage{amsmath}
\usepackage{amsfonts}
\usepackage{amssymb}
\usepackage{booktabs}
\usepackage{geometry}
\usepackage{enumitem}
\usepackage{parskip}
\geometry{margin=2cm}

\title{ZTBSK (2026L), 3. Zarządzanie sieciami komputerowymi}
\author{Piotr Gugnowski, wydział EiTI, nr indeksu 320690}
\date{25 Maj 2026}

\begin{document}
\maketitle

\section*{Pytania i odpowiedzi}

% -----------------------------------------------------------------------
\subsection*{1. Opisz podstawowe obszary działania systemów zarządzania. Jakie
obszary zarządzania są istotne z punktu widzenia administratora: małej sieci,
sieci kampusowej, operatora?}

Standard ISO/IEC~7498-4 definiuje pięć podstawowych obszarów działania systemów
zarządzania, tworzących razem model FCAPS:

\begin{description}
  \item[Zarządzanie obsługą błędów (\emph{Fault Management})]
        Obejmuje wykrywanie, izolowanie i usuwanie niesprawności urządzeń
        i łączy sieciowych w trakcie ich działania. Celem jest jak najszybsze
        przywrócenie sprawności sieci lub odpowiednich jej fragmentów.
  \item[Zarządzanie konfiguracją (\emph{Configuration Management})]
        Pozwala na sprawdzenie aktualnej konfiguracji każdego elementu sieci
        (routerów, przełączników, interfejsów) oraz dokonywanie w niej odpowiednich
        zmian -- zarówno jednorazowych, jak i masowych.
  \item[Rozliczenia (\emph{Accounting Management})]
        Pojęcie rzadko pojawiające się w kontekście sieci TCP/IP, lecz istotne
        dla operatorów. Oznacza monitorowanie kosztów połączeń i wykorzystania
        zasobów sieciowych przez poszczególnych użytkowników lub grupy.
  \item[Zarządzanie wydajnością (\emph{Performance Management})]
        Polega przede wszystkim na monitorowaniu ruchu w sieci w celu określenia
        ewentualnych "wąskich gardeł", pomiarze opóźnień i obciążenia urządzeń.
  \item[Zarządzanie bezpieczeństwem (\emph{Security Management})]
        Obejmuje kontrolę i autoryzację dostępu do zasobów sieciowych, szyfrowanie
        danych oraz alarmowanie o wszelkich przypadkach naruszenia zasad
        bezpieczeństwa.
\end{description}

\textbf{Istotność obszarów z perspektywy różnych administratorów:}

\textbf{Administrator małej sieci.}
Najistotniejsze są zarządzanie konfiguracją (ręczna konfiguracja niewielkiej
liczby urządzeń przez CLI lub interfejs WWW) oraz zarządzanie obsługą błędów
(proste monitorowanie dostępności narzędziami takimi jak MRTG czy \texttt{ping}).
Rozliczenia i zaawansowane zarządzanie wydajnością mają tu mniejsze znaczenie.

\textbf{Administrator sieci kampusowej.}
W sieci obejmującej setki lub tysiące urządzeń kluczowe są cztery obszary
FCAPS -- rozliczenia mają tu marginalne znaczenie, gdyż sieć kampusowa
zazwyczaj nie rozlicza ruchu per użytkownik. Szczególny nacisk kładzie się na:
zarządzanie konfiguracją (zautomatyzowana konfiguracja wielu urządzeń
jednocześnie), zarządzanie obsługą błędów (szybka lokalizacja awarii
w rozległej infrastrukturze), zarządzanie wydajnością (analiza ruchu
w segmentach sieci z użyciem RMON, wykrywanie wąskich gardeł, planowanie
rozbudowy) oraz zarządzanie bezpieczeństwem (ochrona zasobów dostępnych dla
dużej liczby użytkowników).

\textbf{Operator sieci (ISP/telekomunikacyjny).}
Wszystkie obszary FCAPS są krytyczne. Dodatkowo kluczowe stają się rozliczenia
(mierzenie ruchu per klient, per usługa, weryfikacja SLA) oraz wdrożenia
złożonych platform zarządzania klasy enterprise. W dużych sieciach
telekomunikacyjnych stosuje się standard TMN (\emph{Telecommunications
Management Network}, ITU-T), który oferuje znacznie większe możliwości
sterowania i zarządzania niż SNMP.

% -----------------------------------------------------------------------
\subsection*{2. W jaki sposób może odbywać się konfiguracja urządzeń sieciowych?
W jakich sytuacjach które metody wydają się być najwłaściwsze?}

\textbf{Dostępne metody:}

\begin{itemize}
  \item \textbf{Dostęp terminalowy (port szeregowy RS-232 / USB).}
        Pierwotna i nadal powszechna metoda. Urządzenie podłączane jest
        bezpośrednio i nie wymaga działającej sieci IP.
  \item \textbf{Telnet i SSH.}
        Pozwalają na połączenie przez sieć. Telnet przesyła dane jawnym tekstem,
        natomiast SSH szyfruje całą komunikację i jest zdecydowanie preferowany.
  \item \textbf{Interfejs WWW (HTTP/HTTPS).}
        Bardzo upraszcza konfigurację, szczególnie urządzeń przeznaczonych dla
        małych sieci i użytku domowego. Nie wymaga znajomości składni poleceń.
        Dzięki rosnącej mocy obliczeniowej modułów zarządzania możliwe jest
        tworzenie nawet niewielkich systemów zarządzania bezpośrednio w urządzeniu.
  \item \textbf{Linia poleceń -- sieciowe systemy operacyjne.}
        Najważniejsi dostawcy opracowali własne, rozbudowane zestawy komend
        (np.\ Cisco IOS, Juniper JUNOS). Dają one najpełniejszą możliwość
        konfiguracji -- nie wszystkie niuanse konfiguracyjne są możliwe do
        zrealizowania z poziomu interfejsów graficznych.
  \item \textbf{Protokoły zarządzania (SNMP SetRequest, NETCONF).}
        Umożliwiają zdalną, automatyczną konfigurację urządzeń w ramach systemów
        zarządzania, szczególnie istotne przy masowej konfiguracji.
\end{itemize}

\textbf{Dobór metody do sytuacji:}

\begin{itemize}
  \item \textbf{Pierwsze uruchomienie lub awaria sieci} -- port szeregowy jest
        jedyną metodą, gdy sieć IP jeszcze nie działa lub jest niedostępna.
  \item \textbf{Urządzenia domowe i małe biuro} -- interfejs WWW jest
        optymalny: intuicyjny, nie wymaga specjalistycznej wiedzy.
  \item \textbf{Profesjonalna konfiguracja złożonych urządzeń} -- CLI daje
        pełną kontrolę nad wszystkimi parametrami i jest niezbędna, gdy interfejs
        graficzny nie eksponuje określonych opcji konfiguracyjnych.
  \item \textbf{Zdalne zarządzanie przez sieć} -- SSH jest metodą zalecaną
        ze względów bezpieczeństwa (szyfrowanie, uwierzytelnianie kluczem).
        Telnet jest dopuszczalny wyłącznie w izolowanych sieciach zarządzania.
  \item \textbf{Masowa konfiguracja wielu urządzeń} -- protokoły zarządzania
        lub narzędzia do automatyzacji pozwalają skonfigurować setki urządzeń
        jednocześnie, co byłoby niewykonalne metodami manualnymi.
\end{itemize}

% -----------------------------------------------------------------------
\subsection*{3. Omów architekturę systemów zarządzania i rolę poszczególnych
ich elementów.}

Typowa budowa systemu zarządzania jest stosunkowo prosta i składa się
z następujących elementów:

\begin{description}
  \item[Agenci (\emph{agents})]
        Elementy rozmieszczone w zarządzanych urządzeniach, odpowiedzialne za
        monitorowanie i zarządzanie danym obiektem (komputerem, urządzeniem
        sieciowym itp.). Mogą to być fragmenty oprogramowania urządzenia, dedykowane
        moduły sprzętowe lub zupełnie osobne urządzenia (np.\ sondy RMON).
        Przechowują informacje o stanie urządzenia w bazie MIB i odpowiadają na
        zapytania stacji zarządzającej.
  \item[Stacja zarządzająca (\emph{NMS -- Network Management Station})]
        Centralny element systemu. Stale odpytuje agentów o stan urządzeń,
        gromadzi i prezentuje zebrane informacje, generuje alarmy. Komunikacja
        jest zazwyczaj inicjowana przez stację. Agenci inicjują komunikację
        ze stacją w razie sytuacji krytycznych (awaria, przekroczenie
        wcześniej określonych limitów ruchu itp.).
  \item[Urządzenia pośredniczące (\emph{proxy})]
        Stosowane gdy agent i stacja zarządzająca korzystają z różnych protokołów
        komunikacji. Pełnią rolę "tłumacza" pomiędzy nimi, umożliwiając
        zarządzanie urządzeń nieobsługujących natywnie protokołu zarządzania
        stosowanego przez stację.
  \item[Super-agenci (\emph{sub-managers})]
        Stosowane w bardzo rozbudowanych sieciach lub gdy stała bezpośrednia
        łączność pomiędzy agentem a stacją zarządzającą jest zbyt kosztowna
        (np.\ w sieciach rozległych). Zbierają dane od agentów z wyróżnionej
        części sieci, dokonują wstępnej obróbki i agregacji informacji,
        a następnie wysyłają przetworzone dane do centralnej stacji zarządzającej.
\end{description}

Komunikacja między agentami a stacją zarządzającą może odbywać się przez tę
samą sieć, którą zarządzamy (\emph{in-band}), lub przez dedykowaną, odrębną
sieć zarządzania (\emph{out-of-band}). W sieciach telekomunikacyjnych często
stosuje się to drugie rozwiązanie -- ruch zarządzania przesyłany jest wprawdzie
przez tę samą sieć fizyczną, ale wyodrębnionym kanałem, aby nie wpływał na
transmisję danych.

% -----------------------------------------------------------------------
\subsection*{4. Jakie polecenia protokołu SNMP mogą być wykorzystane do
monitorowania stanu sieci?}

Protokół SNMP (\emph{Simple Network Management Protocol}) w swoim podstawowym
modelu posługuje się pięcioma typami komunikatów PDU
(\emph{Protocol Data Unit}):

\begin{description}
  \item[\texttt{GetRequest}]
        Wysyłany przez stację zarządzającą do agenta. Żąda wartości konkretnego,
        pojedynczego elementu bazy danych agenta (jednego obiektu MIB).
  \item[\texttt{GetNextRequest}]
        Wysyłany przez stację zarządzającą. Żąda wartości \emph{następnego}
        elementu bazy danych względem wskazanego. Umożliwia sekwencyjne
        przeglądanie hierarchicznej bazy MIB bez znajomości pełnych identyfikatorów
        obiektów.
  \item[\texttt{SetRequest}]
        Wysyłany przez stację zarządzającą do agenta. Nakazuje wstawienie
        określonej wartości do wskazanego elementu bazy MIB. Zmiana wartości
        w bazie powinna pociągać za sobą odpowiednią zmianę stanu urządzenia.
  \item[\texttt{GetResponse}]
        Wysyłany przez agenta w odpowiedzi na każde z powyższych poleceń.
        Zawiera żądane wartości lub informację o błędzie.
  \item[\texttt{Trap}]
        Jedyny komunikat inicjowany przez \emph{agenta}. Agent wysyła go
        do stacji zarządzania w przypadku nagłego zdarzenia (awaria, przekroczenie
        progu). Zawiera informację o typie zdarzenia, czasie jego powstania
        i ewentualnie dalsze dane o stanie urządzenia.
\end{description}

W wersji SNMPv2 dodano dwa nowe typy: \texttt{GetBulkRequest} (przesłanie
większej liczby danych jednym zapytaniem) oraz \texttt{InformRequest} (wymiana
informacji zarządczych między stacjami zarządzającymi, z potwierdzeniem
odbioru). SNMPv3 nie wprowadził nowych PDU, lecz dodał warstwę
bezpieczeństwa USM (\emph{User Security Model}) -- szyfrowanie (DES),
uwierzytelnianie (MD5, SHA) i zdefiniowane prawa dostępu do danych.

% -----------------------------------------------------------------------
\subsection*{5. Jakiego rodzaju informacje i w jaki sposób są przechowywane
w bazie MIB?}

MIB (\emph{Management Information Base}) to hierarchiczna baza danych
przechowywana przez agenta SNMP. Jej struktura ma postać drzewa:

\begin{itemize}
  \item \textbf{Drzewo obiektów.} Węzły drzewa definiują relacje logiczne
        między zarządzanymi obiektami. Liście drzewa definiują same obiekty
        (właściwości urządzenia, sposób jego działania, dane statystyczne itp.).
  \item \textbf{Rodzaje przechowywanych danych.} Baza MIB przechowuje wyłącznie
        proste obiekty: dane \textbf{skalarne} (pojedyncze wartości) lub
        \textbf{dwuwymiarowe tablice} danych skalarnych. Protokół SNMP obsługuje
        wyłącznie przesyłanie danych skalarnych, więc tablice muszą być
        przesyłane element po elemencie.
  \item \textbf{Identyfikatory obiektów (OID).} Każdy element bazy ma unikalny
        identyfikator będący ciągiem liczb całkowitych oddzielonych kropkami,
        gdzie każda liczba to numer obiektu na danym stopniu hierarchii. Przykładowo
        tablica połączeń TCP dostępna jest pod adresem \texttt{1.3.6.1.2.1.6.13}.
  \item \textbf{Spójność z urządzeniem.} Każda zmiana stanu zarządzanego
        urządzenia powinna natychmiast znajdować odbicie w bazie MIB agenta
        i odwrotnie -- zmiana wartości w bazie MIB powinna powodować odpowiednią
        zmianę stanu urządzenia.
  \item \textbf{Rozszerzenia prywatne.} W hierarchii MIB przewidziano wydzieloną
        gałąź dla wytwórców urządzeń sieciowych -- mogą oni rozszerzać standardową
        bazę MIB o informacje charakterystyczne dla swoich produktów (jest to
        jeden z powodów, dla których nie ma praktycznie systemów zarządzania
        pozwalających na zarządzanie sprzętem różnych producentów -- każdy z nich
        inaczej definiuje parametry swoich urządzeń).
\end{itemize}

% -----------------------------------------------------------------------
\subsection*{6. Jakie grupy informacji zawiera baza danych MIB-II?}

W 1991 roku zaproponowano rozszerzenie definicji bazy danych MIB tworząc
standard MIB-II (RFC~1213). Rozszerza on dane dotąd gromadzone przez bazę MIB
o statystyki dotyczące m.in.\ interfejsów sieciowych, translacji adresów
i protokołów IP, ICMP, UDP, TCP oraz SNMP. MIB-II zawiera następujące grupy
obiektów:

\begin{enumerate}
  \item \textbf{system} -- informacje ogólne o urządzeniu (nazwa, opis,
        czas działania, lokalizacja, dane kontaktowe).
  \item \textbf{interfaces} -- dane o interfejsach sieciowych urządzenia
        (liczba interfejsów, prędkości, stany, statystyki przesłanych pakietów
        i błędów).
  \item \textbf{at} (\emph{address translation}) -- tablice translacji adresów
        sieciowych na fizyczne.
  \item \textbf{ip} -- statystyki protokołu IP (liczba datagramów, tablice
        routingu, błędy fragmentacji).
  \item \textbf{icmp} -- statystyki protokołu ICMP (liczba komunikatów
        poszczególnych typów).
  \item \textbf{tcp} -- statystyki protokołu TCP (liczba połączeń, tablica
        aktywnych połączeń, retransmisje).
  \item \textbf{udp} -- statystyki protokołu UDP.
  \item \textbf{egp} -- statystyki protokołu EGP (\emph{Exterior Gateway Protocol}).
  \item \textbf{transmission} -- dane specyficzne dla medium transmisyjnego.
  \item \textbf{snmp} -- statystyki samego protokołu SNMP (liczba wysłanych
        i odebranych komunikatów poszczególnych typów).
\end{enumerate}

% -----------------------------------------------------------------------
\subsection*{7. Jakie rodzaje informacji o sieci pozwala zbierać standard RMON2?}

RMON (\emph{Remote Network Monitoring}) powstał jako rozszerzenie bazy MIB-II
w celu zapewnienia analizy ruchu w całym segmencie sieci. Pierwotn RMON1
(RFC~1757) dodał dziewięć grup dla sieci Ethernet analizujących ruch na
\textbf{warstwie fizycznej} (statistics, history, alarm, host, hostTopN, matrix,
filter, capture, event) oraz jedną dla sieci Token Ring.

W 1996 roku rozszerzono specyfikację tworząc \textbf{RMON2} (RFC~2021), który
pozwala na analizę ruchu w \textbf{wyższych warstwach modelu sieciowego}.
Nowe grupy RMON2 to:

\begin{description}
  \item[\texttt{protocolDir} (\emph{protocol directory})]
        Katalog wszystkich protokołów sieciowych, które mogą być przez sondę
        interpretowane.
  \item[\texttt{protocolDist} (\emph{protocol distribution})]
        Statystyki ruchu w każdym z protokołów -- umożliwia analizę, jakie
        protokoły dominują w ruchu sieciowym.
  \item[\texttt{addressMap} (\emph{address map})]
        Tablice adresów fizycznych (MAC) i sieciowych (IP) dla wszystkich maszyn
        podłączonych do danego segmentu sieci (mapowanie warstwy 2 na warstwę 3).
  \item[\texttt{nlHost} (\emph{network-layer host})]
        Statystyka ruchu w warstwie sieciowej dla poszczególnych komputerów
        (ruch wysyłany i odbierany per adres IP).
  \item[\texttt{nlMatrix} (\emph{network-layer matrix})]
        Tablica z danymi o ruchu pomiędzy każdymi dwoma komputerami, zebrana
        na podstawie adresów sieciowych -- pozwala stwierdzić, które pary hostów
        generują największy ruch.
  \item[\texttt{alHost} (\emph{application-layer host})]
        Statystyka ruchu dla poszczególnych komputerów zebrana na podstawie ich
        adresów w warstwie aplikacji.
  \item[\texttt{alMatrix} (\emph{application-layer matrix})]
        Tablica z danymi o ruchu pomiędzy każdymi dwoma komputerami, zebrana
        na podstawie adresów warstwy aplikacji.
  \item[\texttt{usrHistory} (\emph{user history collection})]
        Dane historyczne zbierane z innych tabel -- o tym, które wartości będą
        przechowywane, decyduje użytkownik.
  \item[\texttt{probeConfig} (\emph{probe configuration})]
        Standardowe parametry konfiguracyjne sondy RMON.
\end{description}

% -----------------------------------------------------------------------
\subsection*{8. Opisz funkcje znanej platformy zarządzania.}

Ze względu na fakt, że każdy producent urządzeń sieciowych rozszerza standardową
bazę MIB o własne elementy, budowa uniwersalnych platform zarządzania jest
niezwykle trudna -- platform, które byłyby w stanie zarządzać sprzętem
pochodzącym od różnych producentów, istnieje tylko kilka. Przykładem jest
\textbf{HP OpenView}.

\medskip
\textbf{HP OpenView} to wielomodułowa platforma zarządzania siecią, której
główne funkcje obejmują:

\begin{itemize}
  \item \textbf{Wykrywanie topologii sieci.} Platforma automatycznie
        wykrywa urządzenia sieciowe i tworzy graficzną mapę topologii sieci.
  \item \textbf{Monitorowanie dostępności.} Cykliczne sprawdzanie stanu
        urządzeń i usług z generowaniem alarmów w przypadku awarii lub
        przekroczenia progów wydajnościowych.
  \item \textbf{Zarządzanie zdarzeniami i alarmami.} Zbieranie i korelacja
        zdarzeń (w tym komunikatów SNMP Trap) z całej infrastruktury,
        priorytetyzacja alarmów i powiadamianie odpowiednich osób.
  \item \textbf{Monitorowanie wydajności.} Zbieranie danych o obciążeniu
        łączy i urządzeń, generowanie wykresów i raportów wydajności,
        analiza trendów.
  \item \textbf{Zarządzanie konfiguracją.} Dostęp do parametrów konfiguracyjnych
        zarządzanych urządzeń i możliwość ich zmiany.
  \item \textbf{Otwarta architektura.} Możliwość integracji z modułami
        dostarczanymi przez producentów urządzeń sieciowych (NNM SDK),
        co pozwala rozszerzyć platformę o obsługę sprzętu specyficznego
        dla danego wytwórcy.
\end{itemize}

% -----------------------------------------------------------------------
\subsection*{9. Z jakich elementów należałoby zbudować system zarządzania
dla niewielkiej sieci?}

Nawet w małej sieci możliwe jest wykorzystanie możliwości jakie dają zarządzalne
urządzenia sieciowe, przy niewielkim nakładzie kosztów i umiarkowanym nakładzie
pracy:

\begin{enumerate}
  \item \textbf{Zarządzalne urządzenia sieciowe z obsługą SNMP.}
        Routery, przełączniki i punkty dostępowe obsługujące SNMPv3 dla bezpieczeństwa stanowią fundament
        systemu -- umożliwiają zbieranie informacji o stanie sieci.
  \item \textbf{Stacja zarządzająca.}
        Komputer (może być wirtualny) z oprogramowaniem do monitorowania.
        W małej sieci wystarczy jedna stacja.
  \item \textbf{Oprogramowanie monitorowania i wizualizacji.}
        Narzędzia open-source, np.\ MRTG lub Cacti (monitorowanie obciążenia
        interfejsów via SNMP, wykresy historyczne) oraz Nagios lub Zabbix
        (monitorowanie dostępności usług i urządzeń, alarmowanie o awariach).
  \item \textbf{Konfiguracja SNMP na urządzeniach.}
        Na zarządzanych urządzeniach należy skonfigurować tzw.\ community
        string -- hasło dostępu do agenta SNMP. Do monitorowania wystarczy
        tryb tylko do odczytu (ang.\ \emph{read-only}, domyślnie \texttt{public}),
        który pozwala stacji zarządzającej odczytywać dane, ale nie zmieniać
        konfiguracji urządzenia. Należy też wskazać adres stacji zarządzającej
        jako odbiorcę komunikatów Trap oraz -- w miarę możliwości -- zastosować
        SNMPv3 z uwierzytelnianiem i szyfrowaniem zamiast prostego community string.
  \item \textbf{Sieć zarządzania (VLAN zarządzania).}
        W małej sieci komunikacja zarządzania odbywa się zazwyczaj tą samą
        siecią co dane użytkowników. Dobrą praktyką jest wydzielenie osobnej
        sieci VLAN dedykowanej ruchowi zarządzania.
  \item \textbf{Dokumentacja.}
        Mapa sieci, spis urządzeń z adresami IP i typami, procedury reagowania
        na alarmy.
\end{enumerate}

Tak zbudowany system umożliwia wykrywanie awarii, monitorowanie obciążenia
łączy i zarządzanie konfiguracją bez potrzeby ponoszenia kosztów licencyjnych
rozbudowanych platform klasy enterprise.

% -----------------------------------------------------------------------
\subsection*{10. Dlaczego administrator sieci potrzebuje systemu zarządzania?}

Potrzeba posiadania systemu zarządzania wynika z kilku fundamentalnych
przesłanek:

\begin{itemize}
  \item \textbf{Złożoność nowoczesnych sieci.} Nawet stosunkowo mała sieć
        firmowa składa się z dziesiątek urządzeń. Ręczne sprawdzanie stanu
        każdego z nich jest nieefektywne lub wręcz niemożliwe. System
        zarządzania automatyzuje monitorowanie i sygnalizuje problemy.
  \item \textbf{Szybkie wykrywanie i lokalizowanie awarii.} Bez systemu
        zarządzania administrator dowiaduje się o awarii zazwyczaj dopiero
        ze zgłoszenia użytkownika. System wykrywa awarię natychmiast (przez
        cykliczne odpytywanie agentów lub odbiór komunikatów Trap) i wskazuje
        jej lokalizację, co drastycznie skraca czas przywrócenia sprawności
        sieci.
  \item \textbf{Zarządzanie konfiguracją.} System umożliwia zdalne i scentralizowane
        zarządzanie konfiguracją urządzeń, co jest nieocenione przy wdrażaniu
        zmian, aktualizacjach oprogramowania lub reagowaniu na incydenty
        bezpieczeństwa.
  \item \textbf{Monitorowanie wydajności i planowanie rozbudowy.} Zbierane
        statystyki ruchu (przez SNMP i RMON) pozwalają śledzić trendy obciążenia
        łączy, wykrywać wąskie gardła i planować rozbudowę infrastruktury
        zanim pojawią się problemy z wydajnością.
  \item \textbf{Bezpieczeństwo sieci.} System umożliwia monitorowanie prób
        nieautoryzowanego dostępu, alarmowanie o naruszeniach zasad bezpieczeństwa
        oraz szybkie reagowanie na incydenty.
  \item \textbf{Rozliczalność i audyt.} Zebrane dane historyczne umożliwiają
        analizę incydentów po fakcie, generowanie raportów dostępności dla
        kierownictwa oraz weryfikację parametrów SLA.
\end{itemize}

\end{document}
